\section*{Resumen}

La tesis aborda los efectos de escala en la extrapolación modelo--buque, un
problema central de la hidrodinámica naval derivado de la imposibilidad de
igualar simultáneamente los adimensionales entre modelo y buque real. El trabajo
se basa en una metodología combinada CFD--EFD: por un lado, simulaciones CFD
(RANS, con modelos de turbulencia SST y de transición SST-LM) sobre cascos de
pesqueros y hélices tipo Wageningen B4-60, entre otras; por otro, campañas
experimentales de resistencia y de hélices en aguas abiertas realizadas en canal
de ensayos, algunas propias y otras preexistentes en el LabHiNO (previas a la
tesis, o cedidas por una empresa), utilizadas para calibrar y validar los modelos
numéricos. Sobre la resistencia, se estudia en detalle el factor de forma $(1+k)$
de la hipótesis de Hughes: su determinación experimental y numérica combinada, la
influencia de la proa bulbo, la popa sumergida, la relación de aspecto $L/B$ del
casco, la línea de fricción de referencia (ITTC-57 vs.\ líneas numéricas) y el
modelo de turbulencia empleado, evaluando cómo cada uno de estos factores
introduce sesgos sistemáticos en la extrapolación a escala real.

En el capítulo de propulsión, la misma lógica de integración CFD--EFD se aplica al
análisis de los efectos de escala en hélices operando en aguas abiertas, mostrando
las limitaciones del procedimiento estándar ITTC-78 (que asume que toda la
diferencia entre escalas es friccional) frente a la evidencia, respaldada por
simulaciones, de que la componente de presión también escala de forma
significativa, cuantificada mediante la ``fracción de presión''. Se estudia además
el régimen transicional laminar-turbulento a bajos números de Reynolds de modelo,
donde la hipótesis del ITTC-78 falla al no contemplar el comportamiento no monótono
de $K_T/K_Q$, empleando el modelo de transición $\gamma$--$\widetilde{Re}_{\theta t}$
(SST-LM) calibrado y contrastado contra datos experimentales de campañas de aguas
abiertas. En síntesis, es una tesis de hidrodinámica naval con enfoque combinado
numérico--experimental, orientada a mejorar los procedimientos de extrapolación
modelo--buque tanto en resistencia como en propulsión.